\documentclass[11pt]{article}
\usepackage[a4paper,margin=27mm]{geometry}
\usepackage{amsmath,amssymb,amsthm,mathtools,bm,microtype}
\usepackage[colorlinks=true,linkcolor=blue,citecolor=blue,urlcolor=blue]{hyperref}
\newtheorem{theorem}{Theorem}[section]
\newtheorem{lemma}[theorem]{Lemma}
\newtheorem{proposition}[theorem]{Proposition}
\newtheorem{remark}[theorem]{Remark}
\newcommand{\E}{\mathbb E}
\newcommand{\T}{\mathbb T}
\newcommand{\R}{\mathbb R}
\newcommand{\C}{\mathbb C}
\newcommand{\tr}{\operatorname{Tr}}
\newcommand{\eps}{\varepsilon}
\newcommand{\la}[1]{\langle #1\rangle}
\title{Canonical Derivative Currents and Covariance-Normal Energy\protect\\for a Three-Component Fourth-Order Gaussian Field}
\author{Jusang Lee}
\date{Internally audited edition 0.1.2 --- 12 November 2026}
\begin{document}
\maketitle
\begin{abstract}
We study three complex Gaussian fields on a fixed three-dimensional torus,
with a positive matrix-valued fourth-order covariance and a fixed positive
density floor. We construct a rational one-form derivative current by a
common spectral approximation and construct its quadratic energy by exact
conditional covariance subtraction. The current converges in every finite
probability moment with spatial regularity below $-1/2$; the centered energy
converges in $L^2(\Omega;H^{-1-\kappa})$ for each $\kappa>0$.
An explicit finite Gaussian integration-by-parts identity reduces the energy
covariance to three connection families and treats the rational coefficient
without a Taylor or Hermite truncation. For embedded Pauli generators we
derive the exact leading contraction, sharp pointwise compensation slopes,
a subcritical homogeneous instability, and a noncoercive critical boundary.
A determinant estimate supplies a separate bound for deterministic positive
backgrounds. The construction uses a fixed Gaussian reference; construction
of a self-interacting Gibbs measure is outside the conclusions.
\end{abstract}

\section{Introduction and position relative to previous work}
Derivative composites at the regularity considered here require more than
classical multiplication. Paracontrolled calculus supplies the commutator
and paralinearization tools for the first composite \cite{GIP}. Gaussian
contraction and integration by parts supply the second. These mechanisms
are established mathematics; our contribution is their explicit combination
for the fixed rational current below, a regulator comparison for its
conditionally centered energy, and the exact model-dependent stability
boundaries. The Gaussian rough-path constructions of \cite{FV} are useful
comparators, but a one-parameter rough-path theorem is not used to identify
this three-dimensional distribution with a pathwise stochastic integral.
We use the term canonical spectral current throughout.

The analytic BPHZ theorem \cite{CH} is a general stochastic convergence
framework. Its hypotheses are not substituted for the direct estimates
given here. In particular our coefficient is kept as an ordinary smooth
function of the value field, rather than expanded into a globally convergent
power series, and no interacting measure theorem is imported. The fixed
reference, regulator class, and positive floor are part of every assertion.
The companion dependency audit identifies standard inputs and the remaining
model-specific propositions. Bounded literature searches do not certify
priority or replace specialist novelty review.

Related Gaussian gradient-square and analytic Wick-functional results
\cite{CRS,CHRR} concern different covariance orders, rescalings or
correlation limits. They do not directly supply the cross-cutoff Sobolev
convergence for the ordinary rational matrix coefficient in this paper.
The geometric stochastic heat equation \cite{BGHZ} also treats nonlinear
derivative composites, in a different dynamical and spatial setting.
These comparisons rule out a claim to have introduced a general method
for field-dependent derivative squares.

Xu \cite{Xu} proves convergence of rescaled nonlinear Gaussian functionals
to Wick powers under two-sided scalar covariance bounds; those results already
accommodate nonpolynomial nonlinearities. Our object instead retains an
ordinary rational matrix coefficient of a continuous value field and
conditionally centers a quadratic derivative factor, without amplitude
rescaling, on a fixed torus. We prove this mixed composite limit directly
and make no claim to introduce nonpolynomial Gaussian functionals.

\section{Fixed model and conventions}
Let $L=16$, $V=L^3$, $k_n=2\pi n/L$, and write
$\Psi(x)=\sum_{n\in\mathbb Z^3}\widehat\Psi(n)e^{ik_n\cdot x}$.
The field is a centered stationary proper complex Gaussian with
\begin{equation}\label{eq:cov}
 \E[\widehat\Psi(n)\widehat\Psi(m)^\dagger]
 =\frac{2}{V}{\cal A}(k_n)^{-1}\mathbf1_{n=m},\qquad
 \E[\widehat\Psi(n)\widehat\Psi(m)^T]=0,
\end{equation}
where
\begin{align}
 \mathcal A(k)&=(r+Z|k|^2+Y|k|^4)I_3+M,\label{eq:symbol}\\
 M&=\operatorname{diag}(0,0.03,0.07)
       +0.15(I_3-\bm1\bm1^T/3),\nonumber\\
 (r,Z,Y)&=(0.4740336473,-0.9252754126,1).\nonumber
\end{align}
These decimals are exact model inputs. Since $M\ge0$ and
$r-Z^2/(4Y)>0$, the symbol is positive. It is comparable to
$(1+|k|^2)^2I_3$. Realification means
$X=(\operatorname{Re}\Psi,\operatorname{Im}\Psi)\in\R^6$;
using $e_n(x)=V^{-1/2}e^{ik_n\cdot x}$, its orthonormal-mode covariance is $\operatorname{realify}({\cal A}^{-1})$.
The factor two in \eqref{eq:cov} is therefore essential.

Let $T_1,T_2,T_3$ be the Pauli matrices embedded in the first two coordinates
of $\C^3$ and let
$S_A=\left(\begin{smallmatrix}\Re T_A&-\Im T_A\\\Im T_A&\Re T_A\end{smallmatrix}\right)$.
All $S_A$ are real symmetric. Fix $\eps=10^{-12}$ and set
\begin{equation}\label{eq:oneform}
 \rho=|X|^2,\quad m_A=X^TS_AX,\quad q_A=\frac{m_A}{\rho+\eps},
 \quad p_A=2S_AX,\quad v_A=2(S_A-q_AI)X.
\end{equation}
The one-form $v_A\cdot dX=dm_A-q_A\,d\rho=(\rho+\eps)dq_A$
is smooth with polynomially growing derivatives at this fixed floor.
Define the positive semidefinite coefficient matrix
\begin{equation}\label{eq:B}
 B(X)=\sum_{A=1}^3\{a p_Ap_A^T+b(p_Av_A^T+v_Ap_A^T)+c v_Av_A^T\}.
\end{equation}
Here $a=0.2(0.3)^2/d$, $b=0.1(0.3)(0.25)/d$,
$c=0.15(0.25)^2/d$, and $d=2^2+10^{-12}$.
Thus $a,c>0$ and $ac-b^2>0$. Rounded values of $a,b,c$ are
$0.0045,0.001875,0.00234375$; calculations retain $d$ exactly.
The matrix need not be positive definite: the common-phase direction
$i\Psi$ is in its kernel, and it vanishes on a pure third-component field.

An admitted approximation has one real even scalar multiplier $R_N(n)$
on every real component, with $\sup_{N,n}|R_N(n)|<\infty$,
$R_N(n)\to1$ for each $n$, and finite support or sufficient rapid decay
to make $X_N=R_NX$ smooth. All nonlinear products are actual products
of these smooth fields, with exact Fourier convolution. This includes
sharp cubes, sharp balls and scaled smooth cutoffs. No approximation rate
is imposed on a general admitted sequence.

\section{Canonical spectral derivative current}
Fix an even smooth dyadic partition, use $f\prec g$ for its low--high
paraproduct and $f\circ g$ for its comparable-frequency remainder.
Write $\mathcal C^\alpha=B^\alpha_{\infty,\infty}$.

\begin{lemma}[Dyadic reconstruction]\label{lem:reconstruct}
Let $1/3<\alpha<1/2$, $x_N:\T_L^3\to\R^d$ with $d=6$ here, and
$f:\R^d\to(\R^d)^*$. Fiber indices below range from $1$ to $d$.
Suppose smooth approximations $x_N$ converge in
$\mathcal C^\alpha$ and their resonances
$x_N^b\circ\partial_i x_N^a$ converge in $\mathcal C^{2\alpha-1}$.
For every smooth one-form $f$ with polynomially growing derivatives,
$f_a(x_N)\partial_i x_N^a$ has a limit in $\mathcal C^{\alpha-1}$.
The reconstruction is locally continuous in these enhanced inputs, with
polynomial bounds on bounded $x$-ranges for the one-forms \eqref{eq:oneform}.
\end{lemma}
\begin{proof}
The standard dyadic estimates are
$\|u\prec v\|_\beta\lesssim\|u\|_\infty\|v\|_\beta$ and
$\|u\circ v\|_{\beta+\gamma}\lesssim\|u\|_\beta\|v\|_\gamma$
when $\beta+\gamma>0$. Also
$\|u\prec v\|_{\beta+\gamma}\lesssim\|u\|_\beta\|v\|_\gamma$
for $\beta<0$. Taylor expansion on each low-frequency block yields
$f_a(x)=\partial_bf_a(x)\prec x^b+R_a(x)$, with
$R_a(x)\in\mathcal C^{2\alpha}$. The problematic resonance becomes
\begin{align*}
 f_a(x)\circ\partial_i x^a
 ={}&\partial_b f_a(x)(x^b\circ\partial_i x^a)\\
 &+\mathcal C(\partial_bf_a(x),x^b,\partial_i x^a)
       +R_a(x)\circ\partial_i x^a,
\end{align*}
where $\mathcal C(u,v,w)=(u\prec v)\circ w-u(v\circ w)$.
Subtracting $u$ at the high-frequency spatial scale in each block gives a
factor $2^{-j\alpha}$; the remaining two blocks cost
$2^{-j\alpha}2^{-j(\alpha-1)}$. Summation of comparable blocks and their
lower output frequencies gives the commutator in
$\mathcal C^{3\alpha-1}$ when $2\alpha-1<0<3\alpha-1$.
This is the commutator estimate of \cite[Lemma 2.4]{GIP}; its block proof
also gives continuity for differences. The Taylor remainder times
$\partial_i x$ has the same positive regularity. Multiplication of
$\partial_bf_a(x)\in\mathcal C^\alpha$ with the supplied resonance is
allowed since $3\alpha-1>0$. The other paraproducts have regularity at
least $\alpha-1$. Localization to a bounded range of $x$ proves the
claim for polynomially growing derivatives; the fixed rational denominator
makes the localized constants polynomial in that range.
\end{proof}

\begin{theorem}[Canonical current]\label{thm:current}
For every finite $p\ge1$ and $\kappa>0$,
$K_{A,i,N}=v_A(X_N)\cdot\partial_iX_N$ converges in
$L^p(\Omega;\mathcal C^{-1/2-\kappa})$. All admitted approximations
give the same limit. There is an almost surely convergent extracted
subsequence. No local divergent current subtraction is required.
\end{theorem}
\begin{proof}
Write $C^{\rm val}$ for the value-covariance kernel.
The covariance bound $|\widehat C^{\rm val}(n)|\lesssim\la n^{-4}$ gives
$\E|\Delta_j X(x)|^2\lesssim2^{-j}$.
Gaussian moments, Bernstein's inequality with a sufficiently large finite
spatial integrability exponent, and summation give convergence of $X_N$
in every finite probability moment in $\mathcal C^\alpha$, $\alpha<1/2$.
Low modes converge by the multiplier assumption and the high dyadic blocks
are uniformly summable in each slightly weaker norm.

The resonance $\mathbb X_{i,N}^{ba}=X_N^b\circ\partial_iX_N^a$ has zero
expectation by evenness. Wick pairing and the resonant frequency support give
\begin{equation}\label{eq:resvariance}
 \E|\widehat{\mathbb X}_{i,N}^{ba}(n)|^2
 \lesssim\sum_{u+v=n,\ |u|\asymp|v|}
       \la u^{-4}\la v^{-2}\lesssim\la n^{-3}.
\end{equation}
For bounded $|u|\lesssim|n|$ there are $O(\la n^3)$ comparable terms;
the remaining sum is bounded by $\sum_{|u|\gtrsim|n|}\la u^{-6}$.
For each output mode, dominated convergence gives the Cauchy property
for two arbitrary admitted multiplier sequences coupled to the same field.
Second-chaos moment bounds and the same dyadic argument give convergence
in $L^p(\Omega;\mathcal C^{-\eta})$ for every $\eta>0$.
Choose $\alpha\in(1/3,1/2)$ and $\eta<1-2\alpha$.
Lemma~\ref{lem:reconstruct}, polynomial moment localization and a slightly
larger probability exponent yield the result. Choose $\alpha$ sufficiently
close to $1/2$ for the desired $\kappa$. A summable-probability subsequence
follows from convergence in $L^p$. Evenness gives
$\E[X_N^a\partial_iX_N^b]=0$ at every cutoff.
\end{proof}

\begin{remark}[Quantifiers and approximation rates]
The multiplier assumptions alone give no universal rate in $N$.
Multiplying a sharp cutoff by $1-1/\log(N+2)$ preserves them but creates
a slow error in fixed low modes. Quantitative tail estimates for normalized
standard cutoffs cannot be extended to this general class. Almost sure
convergence above is for an extracted subsequence; a full prescribed dyadic
sequence requires an additional summable error estimate.
\end{remark}

The one-form is not generally exact:
$d(v_A\cdot dX)=(\rho+\eps)^{-1}d\rho\wedge dm_A$.
For $m=x^2-y^2$ it equals $-8xy(x^2+y^2+\eps)^{-1}dx\wedge dy$.
Thus a deterministic chain rule alone does not construct this current.
The declared common regulator class is sufficient for the result. Auxiliary
correlated-field anomaly examples do not establish necessity of that class
for the particular diagonal-leading covariance \eqref{eq:cov}.

\section{Covariance-normal energy}
Put $Y_{N,i}=\partial_iX_N$ and
$\Gamma_{N,i}=\E[Y_{N,i}Y_{N,i}^T]$. Define the smooth-cutoff currents
\[
 J_{A,i,N}=p_A(X_N)\cdot Y_{N,i}=\partial_i m_A(X_N),\qquad
 K_{A,i,N}=v_A(X_N)\cdot Y_{N,i}.
\]
Here $J^2_N,JK_N,K^2_N$ mean respectively
$\sum_{A,i}J_{A,i,N}^2$, $\sum_{A,i}J_{A,i,N}K_{A,i,N}$ and
$\sum_{A,i}K_{A,i,N}^2$. Thus
$e_N=(a/2)J^2_N+bJK_N+(c/2)K^2_N$; separate centering below means
subtracting the conditional expectation of each such smooth-cutoff family
given $X_N(x)$, not multiplying limiting distributions.
The exact density is
\begin{equation}\label{eq:energy}
 e_N=\tfrac12\sum_iY_{N,i}^TB(X_N)Y_{N,i},\quad
 C_N(X)=\tfrac12\sum_i\tr[B(X)\Gamma_{N,i}],\quad E_N=e_N-C_N(X_N).
\end{equation}
Independence of the jointly Gaussian vectors $X_N(x)$ and $Y_{N,i}(x)$
implies $C_N(X)=\E[e_N(x)\mid X_N(x)=X]$. In complex notation
\begin{equation}\label{eq:complexct}
 C_N(\Psi)=\sum_{i,A}\{aU_A^\dagger D_{N,i}U_A
 +2b\Re(U_A^\dagger D_{N,i}V_A)+cV_A^\dagger D_{N,i}V_A\},
\end{equation}
where $U_A=T_A\Psi$, $V_A=(T_A-q_AI)\Psi$,
$D_{N,i}=\E[\partial_i\Psi_N\partial_i\Psi_N^\dagger]$, and
$\Gamma_{N,i}=\tfrac12\operatorname{realify}(D_{N,i})$.

\begin{lemma}[Finite Gaussian identity]\label{lem:ibp}
Let $(X,Y,\widetilde X,\widetilde Y)$ be jointly centered Gaussian, with
$\E[XY^T]=\E[\widetilde X\widetilde Y^T]=0$.
Write $P_{a\beta}=\E[Y_a\widetilde X_\beta]$,
$R_{c\alpha}=\E[\widetilde Y_cX_\alpha]$,
$Q_{ac}=\E[Y_a\widetilde Y_c]$ and let $\Gamma,\widetilde\Gamma$
be the two $Y$ covariances. Suppose the smooth coefficients $A,D$ and all
their derivatives through order two have polynomial growth. Then
the expectation of
$A_{ab}(X)(Y_aY_b-\Gamma_{ab})D_{cd}(\widetilde X)
(\widetilde Y_c\widetilde Y_d-\widetilde\Gamma_{cd})$ equals
\begin{align}\label{eq:ibp}
 &(Q_{ac}Q_{bd}+Q_{ad}Q_{bc})\E[A_{ab}D_{cd}]\nonumber\\
 &+\{P_{a\beta}(Q_{bc}R_{d\alpha}+Q_{bd}R_{c\alpha})
       +P_{b\beta}(Q_{ac}R_{d\alpha}+Q_{ad}R_{c\alpha})\}
       \E[\partial_\alpha A_{ab}\partial_\beta D_{cd}]\nonumber\\
 &+P_{a\beta}P_{b\gamma}R_{c\alpha}R_{d\delta}
       \E[\partial_{\alpha\delta}A_{ab}\partial_{\beta\gamma}D_{cd}].
\end{align}
Repeated indices are summed; coefficients on the right are evaluated at
$X$ and $\widetilde X$, respectively.
\end{lemma}
\begin{proof}
Apply $\E[Z_jF(Z)]=\sum_l\operatorname{Cov}(Z_j,Z_l)\E[\partial_lF(Z)]$
to the four displayed derivative variables. Centering cancels same-vertex
$Y$ pairs and the mixed same-vertex covariances are zero. Each remaining
leg pairs with an opposite $Y$ or differentiates the opposite coefficient.
There are two pairings with two $Q$ bridges, four choices with one $Q$ and
one $P$ and $R$, and one coefficient-derivative term with two $P$ and $R$.
These are exactly \eqref{eq:ibp}. Polynomial growth justifies Gaussian
integration by parts by smooth truncation. The identity differentiates
either coefficient at most twice.
\end{proof}

\begin{theorem}[Energy convergence]\label{thm:energy}
For every $\kappa>0$, $E_N$ converges in
$L^2(\Omega;H^{-1-\kappa}(\T_L^3))$, independently of the admitted
multiplier sequence. The separately centered $J^2,JK,K^2$ converge jointly
in that space. Also $\int E_N\,dx$ converges in $L^2(\Omega)$ with mean zero.
\end{theorem}
\begin{proof}
Use common underlying Fourier modes for all approximations. Cross-cutoff
value covariances have Fourier coefficients bounded by $C\la n^{-4}$;
one-derivative and two-derivative cross-covariances $P,R,Q$ are bounded by
$C\la n^{-3}$ and $C\la n^{-2}$. Value covariances converge uniformly
in space since $\sum_n\la n^{-4}<\infty$ in dimension three.
Expectations of products of coefficient derivatives through order two
therefore converge uniformly: localize the jointly Gaussian values to a
fixed ball, use continuity in the finite covariance matrix there, and remove
the localization by polynomial moment bounds uniform in the cutoffs.

Let $a_t(n)=\la n^{-t}$ and let $G_s$ be the periodic positive Bessel kernel
with Fourier coefficients $\la n^{-2s}$, $s>1$.
Its heat-kernel representation proves positivity. Splitting convolution
sums into the smaller-than-half-output region and dyadic complementary
annuli gives
\begin{equation}\label{eq:conv}
 a_2*a_2\lesssim a_1,\qquad
 a_3*a_3*a_3*a_3\lesssim_\eta a_{3-\eta}\quad(\eta>0).
\end{equation}
Indeed convolution with the dimension-critical $a_3$ adds at most a
logarithm; three such logarithms are bounded by $C_\eta\la n^\eta$.
The first estimate follows directly from the two regions
$|u|\le|n|/2$ and $|n-u|\le|n|/2$ and the remaining annuli.

Expand $\int G_s|Q_{NM}|^2$ in Fourier series and bound each summand
absolutely by $C\la n^{-2s}(a_2*a_2)(n)$.
This majorant is summable exactly for $s>1$. Dominated convergence on the
double sum proves strong convergence of $Q_{NM}$ in $L^2(G_s\,dx)$ as
both cutoffs tend to infinity, not just a uniform norm bound.
The fourfold version and the second estimate of \eqref{eq:conv} prove strong
convergence of $P_{NM},R_{NM}$ in $L^4(G_s\,dx)$.
Consequently $Q^2$, $QPR$ and $P^2R^2$ converge strongly in $L^1(G_s\,dx)$,
the middle assertion following by Holder's inequality.

Apply Lemma~\ref{lem:ibp} with $A,D$ equal to the coefficient matrices in
\eqref{eq:B}. The preceding strong convergence and uniform coefficient
convergence make each cross-cutoff energy covariance converge against $G_s$.
Stationarity and Fourier duality identify its integral, up to the fixed
normalization $V$, with the squared $H^{-s}$ probability norm.
Polarization of the four covariances for $E_N-E_M$ proves the Cauchy
property. The same argument works for each summand of $B$ and for two
different regulator sequences. Pairing with the constant function gives
the integrated assertion. Conditional centering proves zero mean.
\end{proof}

The argument excludes $H^{-1}$; it does not prove that this endpoint fails
for the full rational coefficient. Centering is exact for each regulator;
subtracting only the leading divergent term is a different prescription.

\section{Explicit contraction and stability boundaries}
For sharp cubes $|n|_\infty\le N$, let
\begin{equation}\label{eq:delta}
 \delta=\frac{I_\square}{6\pi^2YL},\qquad
 I_\square=24\int_0^1\!\int_0^1\frac{du\,dv}{1+u^2+v^2}.
\end{equation}
Riemann summation of the homogeneous leading symbol gives
$D_{N,i}/N\to\delta I_3$; lower-order symbol corrections are summable
after differentiation. To justify the Riemann sum near zero, remove a
fixed small cube there, apply ordinary summation elsewhere, and bound
its contribution by $C\eta$ uniformly after rescaling.
The face formula for $I_\square$ follows from
$\operatorname{div}(u/|u|^2)=|u|^{-2}$ and the divergence theorem.
Numerically $\delta\simeq0.01619898645075695$.

Set $s=|\Psi_1|^2+|\Psi_2|^2$ and $\rho=|\Psi|^2$.
The Pauli identities $\sum_A m_A^2=s^2$,
$\sum_A\|T_A\Psi\|^2=3s$ imply
\begin{equation}\label{eq:W}
 \frac{C_N(\Psi)}{N}\longrightarrow\delta W_\eps(\Psi),\qquad
 W_\eps=hs-\frac{6bs^2}{\rho+\eps}
             -\frac{3cs^2(\rho+2\eps)}{(\rho+\eps)^2},
 \quad h=9(a+2b+c).
\end{equation}
The $s^2$ terms follow by expanding
$\sum_A\Re(U_A^\dagger V_A)=3s-s^2/(\rho+\eps)$ and
$\sum_A\|V_A\|^2=3s-s^2(\rho+2\eps)/(\rho+\eps)^2$.

\begin{proposition}[Sharp pointwise slopes]\label{prop:slopes}
For the positive production coefficients, put $\ell=9a+12b+6c$.
For all fields,
\begin{equation}\label{eq:slopes}
 \ell s\le W_\eps\le hs,\qquad
 \inf_{\Psi:s(\Psi)>0}W_\eps/s=\ell,\quad
 \sup_{\Psi:s(\Psi)>0}W_\eps/s=h.
\end{equation}
The upper bound $ms\ge W_\eps$ holds for every field if and only if $m\ge h$.
At $m=h$ the nonnegative residual has no positive lower bound proportional
to $s$, even when $s$ is fixed and nonzero.
\end{proposition}
\begin{proof}
The upper residual is the sum of the last two nonnegative terms in
\eqref{eq:W}. For the lower bound write
\begin{align*}
 W_\eps-\ell s={}&6bs\left(1-\frac{s}{\rho+\eps}\right)\\
 &+3cs\left(1-\frac{s(\rho+2\eps)}{(\rho+\eps)^2}\right)\ge0,
\end{align*}
since $0\le s\le\rho$ and $\rho(\rho+2\eps)\le(\rho+\eps)^2$.
Taking $\rho=s\to\infty$ proves sharpness of $\ell$.
Taking fixed $s>0$ and $\rho\to\infty$ proves sharpness of $h$ and the
failure of a positive uniform proportional bound for $hs-W_\eps$.
\end{proof}

\begin{proposition}[Subcritical homogeneous instability]\label{prop:instability}
Let $t\to\infty$, $m<h$ and $\gamma=1.62$. Let $G_{\rm lo}$ be independent
of $t$, with $|G_{\rm lo}(u)|\le C(1+|u|^4)$.
The homogeneous energy density
$\gamma\rho^3/6+G_{\rm lo}(u)-tW_\eps(u)+tms(u)$
is at most $-c_m t^{3/2}+O(t)$ along suitable fields, for some $c_m>0$.
The fixed production instance is
$G_{\rm lo}(u)=\tfrac12u^\dagger(rI_3+M)u+(\lambda/4)\rho^2$,
with $\lambda=-0.43$ and $\rho=|u|^2$.
\end{proposition}
\begin{proof}
Choose a fixed $x\in(0,1)$ small enough that
$A_x=(h-m)x-(6b+3c)x^2>0$ and take the constant field
$u=(\sqrt{x\rho},0,\sqrt{(1-x)\rho})$, so $s=x\rho$.
At large $\rho$, $W_\eps-ms=A_x\rho+O(1)$.
All derivative currents vanish for this homogeneous field. The local
energy is $\gamma\rho^3/6-tA_x\rho+O(\rho^2+t)$.
Choose $\rho=(2tA_x/\gamma)^{1/2}$. The leading value is
$-2^{3/2}(tA_x)^{3/2}/(3\sqrt\gamma)$, and the error is $O(t)$.
\end{proof}

Thus the minimal leading compensation slope for a family energy
$\tfrac12\Delta\mu_Ns$ is $\Delta\mu_N/N\to2h\delta$
when only leading pointwise cancellation is considered. We make no claim
that a remainder merely denoted $O(1)$ repairs the exact finite-cutoff
counterterm uniformly in all field directions. The threshold is a local
compensation statement, not a full-field coercivity theorem.

\section{A deterministic-background exponential bound}
This separate result gives a useful positive stability estimate. Let $b(x)$
be a deterministic bounded positive semidefinite $18\times18$ matrix on the
bundle of all six-real spatial derivatives. Write $\|b\|_\infty$ for its
essential supremum operator norm. Let $C$ be the Gaussian covariance operator
and $L_b=\nabla^*b\nabla$ in quadratic-form sense. Set
$T_b=C^{1/2}L_bC^{1/2}$, $C_N^{\rm op}=R_NCR_N$ and
$T_{b,N}=R_NT_bR_N$ for the same common multiplier.
The centered quadratic form is $Q_{b,N}=\tfrac12(\la{X_N,L_bX_N}
-\tr(C_N^{\rm op}L_b))$.

\begin{proposition}[Explicit frozen-background bound]\label{prop:det}
Let $M_R=\sup|R_N|$ and
$c_0=(r+Y-\sqrt{(r-Y)^2+Z^2})/4>0$.
For every $p\ge0$ the forms converge in $L^2(\Omega)$ and
\begin{align}\label{eq:det}
 0\le\log\E e^{-pQ_{b,N}}
 &\le\frac{p^2}{4}\|T_{b,N}\|_{\mathrm{HS}}^2\nonumber\\
 &\le\frac{3p^2M_R^4\|b\|_\infty^2}{2c_0^2}
       \left(\frac{L}{2\pi}\right)^4
       \left(4\pi^2+\frac{\pi^4}{45}\right).
\end{align}
The partition factors converge to $\det_2(I+pT_b)^{-1/2}$.
\end{proposition}
\begin{proof}
The smallest eigenvalue of
$\left(\begin{smallmatrix}r&Z/2\\Z/2&Y\end{smallmatrix}\right)$,
and $1+t^2\ge(1+t)^2/2$, give
${\cal A}(k)\ge c_0(1+|k|^2)^2I$.
Quadratic-form order gives
$0\le T_{b,N}\le M_R^2\|b\|_\infty C^{1/2}(-\Delta)C^{1/2}$.
Indeed $b\le\|b\|_\infty I$ first gives a bound by
$\|b\|_\infty R_N C^{1/2}(-\Delta)C^{1/2}R_N$.
Only the translation-invariant middle operator must commute with $R_N$;
no commutation between $b(x)$ and $R_N$ is assumed.
Eigenvalue comparison bounds the squared Hilbert--Schmidt norm by
$6M_R^4\|b\|_\infty^2c_0^{-2}\sum_{n\ne0}|k_n|^{-4}$.
Counting the $24j^2+2$ integer points on the shell $|n|_\infty=j$
gives $\sum_{n\ne0}|n|^{-4}\le24\zeta(2)+2\zeta(4)
=4\pi^2+\pi^4/45$.
The limiting operator is Hilbert--Schmidt, and strong convergence of bounded
multipliers implies its cutoff conjugations converge in that norm.
The common Gaussian coupling and Wick pairing give
$\E|Q_{b,N}-Q_{b,M}|^2=\tfrac12\|T_{b,N}-T_{b,M}\|_{\rm HS}^2$,
which proves the asserted $L^2$ convergence.
Gaussian diagonalization yields
\[
 \log\E e^{-pQ_{b,N}}=\tfrac12\sum_j
       [p\lambda_{j,N}-\log(1+p\lambda_{j,N})].
\]
Use $0\le x-\log(1+x)\le x^2/2$ for $x\ge0$ and Hilbert--Schmidt
convergence. For completeness, along a positive operator segment the
trace derivative of $F_p(T)=\tfrac12\tr[pT-\log(I+pT)]$ is bounded by
$\tfrac{p^2}{2}\|T\|_{\rm HS}$ in the Hilbert--Schmidt dual norm.
Integration along the segment gives continuity in that norm and the
claimed regularized determinant limit.
\end{proof}

For a deterministic bounded $u$, one can take $b=I_3\otimes B(u)$, since
$\|B(u)\|\le(12a+48|b|+48c)|u|^2$.
Replacing $u$ by the random field in this proof would change the Gaussian
diagonalization step and is not justified.

\section{Scope, reproducibility and outstanding review}
The statements are about composites under \eqref{eq:cov}, at fixed $L$ and
fixed $\eps$. The exact phase field $\Psi(x)=e^{ik\cdot x}u$, with
$k\in(2\pi/L)\mathbb Z^3$, has
$J_A=K_A=0$ while $W_\eps(u)>0$ when $s>0$.
Consequently the conditional contraction is not a deterministic rate
function for the bare field law. Likewise $L^2$ convergence of the integrated
centered energy does not imply uniform integrability of its exponential.
The self-coupled negative-exponential estimate, an interacting Gibbs limit,
the A13 constructive program and a physical vacuum are not conclusions.

The verification package separates exact algebra and finite Gaussian moment
checks from the analytic proofs above. It retains original source hashes,
tests the real/complex covariance factor, tests all three integration-by-parts
families with nonconstant coefficients, and rejects altered coefficients and
stronger scope assertions. It also records the slow-cutoff negative control
and the critical third-component escape. Numerical checks are reproducibility
and falsification aids; they are not substitutes for the analytic limits.

This version has undergone internal content review. Independent human
mathematical review and specialist novelty review have not been performed.
The external-review handoff records their disposition separately from file
integrity and code execution. Author declarations, venue choice and actual
transmission require the author's decisions.

\begin{thebibliography}{9}
\footnotesize
\setlength{\itemsep}{2pt}
\bibitem{GIP} M. Gubinelli, P. Imkeller and N. Perkowski,
\emph{Paracontrolled distributions and singular PDEs}, Forum Math. Pi
3 (2015), e6, 75 pp. \url{https://doi.org/10.1017/fmp.2015.2}.
\bibitem{FV} P. Friz and N. Victoir,
\emph{Differential equations driven by Gaussian signals}, Ann. Inst. H.
Poincare Probab. Statist. 46 (2010), 369--413.
\url{https://doi.org/10.1214/09-AIHP202}.
\bibitem{CH} A. Chandra and M. Hairer,
\emph{An analytic BPHZ theorem for regularity structures}, preprint,
arXiv:1612.08138 (2016), revised version v5 (22 January 2018).
\url{https://arxiv.org/abs/1612.08138v5}.
\bibitem{CRS} F. Coppini, W. Ruszel and D. Schuricht,
\emph{Wick theorem for analytic functions of Gaussian fields}, preprint,
arXiv:2507.05131. \url{https://arxiv.org/abs/2507.05131}.
\bibitem{CHRR} A. Cipriani, R. S. Hazra, A. Rapoport and W. Ruszel,
\emph{Properties of the Gradient Squared of the Discrete Gaussian Free Field},
J. Stat. Phys. 190 (2023), 171.
\url{https://doi.org/10.1007/s10955-023-03187-3}.
\bibitem{BGHZ} Y. Bruned, F. Gabriel, M. Hairer and L. Zambotti,
\emph{Geometric stochastic heat equations},
\url{https://arxiv.org/abs/1902.02884}.
\bibitem{Xu} W. Xu,
\emph{Sharp Convergence of Nonlinear Functionals of a Class of Gaussian Random Fields},
Commun. Math. Stat. 6 (2018), 509--532.
\url{https://doi.org/10.1007/s40304-018-0162-9}.
\end{thebibliography}
\end{document}
